\documentclass[12pt]{article}
\usepackage{pst-all}
\usepackage[french]{babel}
\usepackage{tikz}
\usepackage{tkz-tab}
\usepackage[T1]{fontenc}
\usepackage{eurosym}
\usepackage{amssymb}
\usepackage{float}
\usepackage{blkarray}
\usepackage{dsfont}
\usepackage{amsfonts}
\usepackage{amsmath}
\usepackage{array}
\usepackage{graphicx}
\parindent=0cm
\usepackage{geometry}
\usepackage{color}
\usepackage{graphicx}
\geometry{hmargin=1.5cm,vmargin=1.5cm}
\usepackage{amsmath, amssymb, amsfonts, amscd, amsthm}

\usepackage{fancyhdr} 
\pagestyle{fancy}
\usepackage{multirow}
\renewcommand{\headrulewidth}{1pt}
\renewcommand{\footrulewidth}{1pt}
\fancyhead[L]{~}
\fancyhead[R]{~}
\fancyhead[C]{$5^e$}





\usepackage[utf8]{inputenc}
\usepackage{tcolorbox}% charge aussi pgf, verbatim, etoolbox et environ
\tcbuselibrary{theorems}% charge aussi amsmath
\usepackage{xcolor}

\definecolor{propcolor}{rgb}{.804,.064,.464}
\definecolor{defcolor}{rgb}{0,.545,0}
\definecolor{thcolor}{rgb}{0.932,0,0}
\definecolor{remcolor}{rgb}{1,0.5,0}
\definecolor{excolor}{rgb}{.094,.455,.804}

\usepackage{cleveref}

\newtcbtheorem[crefname={definition}{definitions}]{Definition}{Definition}{defstyle}{def}


\usepackage{varwidth} 
\tcbuselibrary{skins}

\newtcbtheorem[use counter from=Definition]{defin}{Définition}
{enhanced,frame empty, theorem name,interior empty,colframe=defcolor,coltitle=defcolor,fonttitle=\bfseries,
	colbacktitle=white,
	borderline={0.5mm}{0mm}{defcolor},
	%borderline={0.5mm}{0mm}{ForestGreen!50!white,dashed},
	attach boxed title to top left={yshift=-2mm, xshift = 2mm},
	boxed title style={boxrule=0.4pt}, varwidth boxed title}{theo}

\newtcbtheorem[use counter from=Definition]{theo}{Théorème}
{enhanced,frame empty,theorem name,interior empty,colframe=thcolor,coltitle=thcolor,fonttitle=\bfseries,
	colbacktitle=white,
	borderline={0.5mm}{0mm}{thcolor},
	%borderline={0.5mm}{0mm}{ForestGreen!50!white,dashed},
	attach boxed title to top left={yshift=-2mm, xshift = 2mm},
	boxed title style={boxrule=0.4pt}, varwidth boxed title}{theo}

\newtcbtheorem[use counter from=Definition]{prop}{Propriété}
{enhanced,frame empty,theorem name,interior empty,colframe=propcolor,coltitle=propcolor,fonttitle=\bfseries,
	colbacktitle=white,
	borderline={0.5mm}{0mm}{propcolor},
	%borderline={0.5mm}{0mm}{ForestGreen!50!white,dashed},
	attach boxed title to top left={yshift=-2mm, xshift = 2mm},
	boxed title style={boxrule=0.4pt}, varwidth boxed title}{theo}


\newtcbtheorem[use counter from=Definition]{ex}{Exemple}
{enhanced,frame empty,theorem name,interior empty,colframe=excolor,coltitle=excolor,fonttitle=\bfseries,
	colbacktitle=white,
	borderline={0.5mm}{0mm}{excolor},
	%borderline={0.5mm}{0mm}{ForestGreen!50!white,dashed},
	attach boxed title to top left={yshift=-2mm, xshift = 2mm},
	boxed title style={boxrule=0.4pt}, varwidth boxed title}{theo}

\newtcbtheorem[use counter from=Definition]{rem}{Remarque}
{enhanced,frame empty,theorem name,interior empty,colframe=remcolor,coltitle=remcolor,fonttitle=\bfseries,
	colbacktitle=white,
	borderline={0.5mm}{0mm}{remcolor},
	%borderline={0.5mm}{0mm}{ForestGreen!50!white,dashed},
	attach boxed title to top left={yshift=-2mm, xshift = 2mm},
	boxed title style={boxrule=0.4pt}, varwidth boxed title}{theo}


\newif\iftrou\troutrue
%mettre true à la place de false pour mettre des pointillés

\newlength{\trouwidth}

\newcommand{\trou}[1]{%
	\iftrou%
	\settowidth{\trouwidth}{#1}%
	\hbox to \trouwidth{\dotfill}%
	\else%
	#1%
	\fi}


\newcommand\pointille[1][1cm]{\makebox[#1]{\dotfill}}


\begin{document}





%%%%%%%%%%%%%%%%%%%%%%%%%%%%%%%%%%%%%%%%%%%%%%%%%%%%%%%%%%%%


\begin{center}
		\rule{\linewidth}{0.5mm} \\[0.4cm]
		{ \huge \bfseries Nombres entiers \\[0.4cm] }
		\rule{\linewidth}{0.5mm} \\
\end{center}



%%%%%%%%%%%%%%%%%%%%%%%%%%%%%%%%%%%%%%%%%%%%%%%%%%%%%%%%%%%%%%%%%%%%%%

\section{Diviseurs d'un nombre entier}


\begin{defin}{}{}
	Un \textbf{entier naturel} est un nombre entier positif ou nul.
\end{defin}

\begin{ex}{}{}
	$1$, $3$, $0$, $6$ $873$ $261$ sont des entiers naturels.
\end{ex}

\begin{defin}{}{}
	Effectuer la \textbf{division euclidienne} d'un nombre entier $a$ par un nombre entier $b$ différent de $0$, c'est trouver deux nombres entiers naturels $q$ et $r$ tels que: \textcolor{olive}{$a$} = \textcolor{blue}{$b$} $\times$ \textcolor{red}{$q$} + \textcolor{orange}{$r$}\\
	\textcolor{olive}{$a$} s'appelle le \textcolor{olive}{dividende}, \textcolor{blue}{$b$} le \textcolor{blue}{diviseur}, \textcolor{red}{$q$} le \textcolor{red}{quotient} et \textcolor{orange}{$r$} le \textcolor{orange}{reste}.
\end{defin} 



\begin{ex}{}{}
	La division euclidienne de $377$ par $12$:
	\begin{center}
		\includegraphics[width=0.35\linewidth]{div1}
	\end{center}
\end{ex}

\begin{defin}{}{}
	Soient $a$ et $b$ deux nombres entiers positifs avec $b \ne 0$:\\
	Lorsque la division euclidienne de $a$ par $b$ donne un reste \textbf{nul}, on dit que :\\
	- $a$ est un \textbf{multiple} de $b$;\\
	- $b$ est un \textbf{diviseur} de $a$;\\
	- $a$ est \textbf{divisible} par $b$.
\end{defin}

\begin{ex}{}{}
	\begin{enumerate}
		\item $20$ est un multiple de \trou{$1$; $2$; $4$; $5$; $10$ et $20$}.
		\item On peut aussi dire que \trou{$1$; $2$; $4$; $5$; $10$ et $20$} sont des diviseurs de $20$.
		\item On dit aussi que $20$ est divisible par  \trou{$1$; $2$; $4$; $5$; $10$ et $20$}.
	\end{enumerate}
\end{ex}

\begin{rem}{}{}
	Lorsqu'on trouve que le reste de la division de est nul, on peut l'écrire sous forme de fraction. 
	\\
	Par exemple $15=3\times 5 +0$, donc \trou{$\dfrac{15}{5}=3$}, et \trou{$\dfrac{15}{3}=5$}.
\end{rem}


\begin{ex}{}{}
	La division euclidienne de $85$ par $17$ donne : \trou{$85 = 17 \times 5 + 0$}, le reste est \trou{nul}. On dit que:\\
	- $85$ est un multiple de $17$ et de $5$;\\
	- $5$ et $17$ sont les diviseurs de $85$;\\
	- $85$ est divisible par $17$ et par $5$.
\end{ex}

\begin{rem}{}{}
	$0$ ne divise aucun nombre.\\
	Tout nombre est un multiple de \trou{$~1~$}.\\
	Tout nombre entier est divisible par \trou{$1$ et par lui-même}.
\end{rem}



%%%%%%%%%%%%%%%%%%%%%%%%%%%%%%%%%%%%%%%%%%%%%%%%%%%%%


\section{Critères de divisibilité}

\begin{prop}{}{}
	\begin{enumerate}
		\item Si un nombre entier a pour chiffre des unités $0$, $2$, $4$, $6$ ou $8$, alors il est \textbf{divisible par $2$}.
		\item Si la somme des chiffres d'un nombre entier est divisible par $3$, alors ce nombre est \textbf{divisible par $3$}.
%		\item Si les deux derniers chiffres d'un nombre entier forment un nombre divisible par $4$, alors ce nombre est \textbf{divisible par $4$}.
		\item Si un nombre entier a pour chiffre des unités $0$ ou $5$, alors il est \textbf{divisible par $5$}.
		\item Si la somme des chiffres d'un nombre entier est divisible par $9$, alors ce nombre est \textbf{divisible par $9$}.
		\item Si un nombre entier a pour chiffre des unités $0$, alors il est \textbf{divisible par $10$}.
	\end{enumerate}
\end{prop}

\begin{ex}{}{}
	\begin{enumerate}
		\item $138$ est divisible par \trou{$2$} car son chiffre des unités est \trou{$8$}, c'est un nombre \trou{pair}.
		\item $186$ est divisible par \trou{$3$} car la somme de ses chiffres est égale à \trou{$15$}, qui est divisible par \trou{$3$}. 
%		\item $5$ $436$ est divisible par \trou{$4$} car $36$ est divisible par \trou{$4$}.
		\item $175$ est divisible par \trou{$5$} car son chiffre des unités est \trou{$5$}.
		\item $936$ est divisible par \trou{$9$} car la somme de ses chiffres est égale à \trou{$18$}, qui est divisible par \trou{$9$}.
		\item $220$ est divisible par \trou{$10$} car son chiffre des unités est \trou{$0$}.
	\end{enumerate}
\end{ex}

%%%%%%%%%%%%%%%%%%%%%%%%%%%%%%%%%%%%%%%%%%%%%%%%%%%%%
\newpage

\section{Nombre premier}

\begin{defin}{}{}
	Un \textbf{nombre premier} est un nombre entier positif qui admet exactement deux diviseurs : $1$ et lui-même.
\end{defin}

\begin{ex}{}{}
\begin{enumerate}
	\item $6$ n'est pas un nombre premier : il admet $2$ et $3$ comme diviseurs.
	\item $7$ est un nombre premier : il n'est divisible que par $1$ et par $7$.
\end{enumerate}
\end{ex}

\begin{rem}{}{}
	Il existe une infinité de nombres premiers.\\
	Le nombre $1$ n’est pas premier car il n’a qu’un seul diviseur.
\end{rem}

\begin{prop}{}{}
Liste des nombres premiers inférieurs à $30$: \trou{$2$; $3$; $5$; $7$; $11$; $13$; $17$; $19$; $23$; $29$}
\end{prop}


\begin{prop}{Décomposition en produit de facteurs premiers}{}
On peut décomposer chaque entier naturel supérieur ou égal $2$ en produits de facteurs premiers.\\
On divise le nombre à décomposer autant de fois que possible par $2$, puis par $3$, par $5$, par $7$, par $11$, etc, en suivant la liste des nombres premiers successifs.\\
\end{prop}

\begin{ex}{}{}
	\begin{enumerate}
		\item 
	$20 = 2 \times 2 \times 5$ est une décomposition du nombre $20$ en produits de facteurs premiers. \\
	En effet, chaque facteur de la décomposition est un nombre premier. 
	\item  $231 =$ \trou{$3 \times 7 \times 11$}
	\item $225 =$ \trou{$3 \times 3 \times 5 \times 5$}
		\end{enumerate}
\end{ex}


%%%%%%%%%%%%%%%%%%%%%%%%%%%%%%%%%%%%%%%%%%%%%%%%%%%%%
\newpage
\section{Applications des nombres premiers aux fractions}


\begin{prop}{}{}
	Pour simplifier une fraction, on décompose son numérateur et son dénominateur en produits de facteurs premiers.
\end{prop}


\begin{ex}{}{}
	Simplifier la fraction $\displaystyle\frac{36}{42}$:\\
	\begin{enumerate}
		\item On décompose $36$ en produit de facteurs premiers : $36 = $ \trou{$2 \times 2 \times 3 \times 3$}
		\item On décompose $42$ en produit de facteurs premiers : $42 = $ \trou{$2 \times 3 \times 7$}
		\item On simplifie : $\displaystyle\frac{36}{42} = \displaystyle\frac{2 \times \textcolor{red}{2 \times 3} \times 3}{\textcolor{red}{2 \times 3} \times 7} = \displaystyle\frac{2 \times 3 }{7} = \displaystyle\frac{6}{7}$
	\end{enumerate}
\end{ex}


\begin{ex}{}{}
	Simplifier la fraction $\displaystyle\frac{153}{85}$:\\
	\begin{enumerate}
		\item On décompose $153$ en produit de facteurs premiers : $153 = $ \trou{$3 \times 3 \times 17$}
		\item On décompose $85$ en produit de facteurs premiers : $85 = $ \trou{$5 \times 17$}
		\item On simplifie : $\displaystyle\frac{153}{85} = \displaystyle\frac{3 \times 3 \times \textcolor{red}{17}}{5 \times \textcolor{red}{17}} = \displaystyle\frac{3 \times 3 }{5} = \displaystyle\frac{9}{5}$
	\end{enumerate}
\end{ex}



%%%%%%%%%%%%%%%%%%%%%%%%%%%%%%%%%%%%%%%%%%%%%%%%%%%%%

\section{Carte mentale}

\begin{center}
		\includegraphics[width=0.99\linewidth]{cmnbentier}
\end{center}



\end{document}
